% GENERATED FILE. Edit canonical lessons, then run npm run generate:books.
% canonical-source-hash: fnv1a64:43e115147fc175e9
% canonical-lessons: TE-C440-alasyam, TE-C440-veci-undu, TE-C440-samayapalana, TE-C440-viramam, TE-C440-tvarapadu

\chapter{Waiting on the Platform}
\label{ch:te-waiting-on-the-platform}
\hlchaptermodality{440}

\begin{chapteropening}
\textbf{By the end of this chapter:} \emph{I can say a delay, to wait, punctuality, a break and to hurry.}

Complete the last lesson of chapter 440: I can say a delay, to wait, punctuality, a break and to hurry.
\end{chapteropening}

\section[\texorpdfstring{\te{ఆలస్యం} (ālasyaṁ)}{ఆలస్యం (ālasyaṁ)}]{\te{ఆలస్యం} (ālasyaṁ) — a delay; lateness}

\label{lesson:TE-C440-alasyam}



\begin{quote}
Before the new one: say the Telugu for a railway coach, then the Telugu for a fare.
\end{quote}

\subsection*{You'll want to know: \te{ఆలస్యం}}
\textbf{\te{ఆలస్యం}} — \emph{ālasyaṁ} — \textquotedblleft{}a delay; lateness\textquotedblright{}.

\begin{quote}
\te{వర్షం} \te{వల్ల} \te{రైలు} \te{రెండు} \te{గంటలు} \te{ఆలస్యం} \te{అయింది}.
\end{quote}

\begin{quote}
\emph{varṣaṁ valla railu reṇḍu gaṇṭalu ālasyaṁ ayindi.}
\end{quote}

\begin{quote}
Because of the rain the train was two hours late.
\end{quote}

\te{ఆలస్యంగా} is late, adverb: \te{ఆలస్యంగా} \te{వచ్చాడు}, he came late. In Sanskrit the word meant laziness; Telugu uses it for lateness.

\begin{cousinweb}
\textbf{\te{ఆలస్యం}} (\emph{ālasyaṁ}) is borrowed from Sanskrit \emph{ālasya}, \textquotedblleft{}sloth, idleness\textquotedblright{} (CDIAL 1371). In Telugu the slowness of the idle became plain lateness.
\end{cousinweb}

\begin{grammarlens}[title={What you've built}]
One more word for getting around.
\end{grammarlens}

\subsection*{Your turn}
\begin{itemize}
\raggedright
  \item \emph{Say it:} \emph{ālasyaṁ}
  \item \emph{Say it:} \emph{ālasyaṁ}, once more
  \item \emph{Say it:} \emph{varṣaṁ valla railu reṇḍu gaṇṭalu ālasyaṁ ayindi.} — Because of the rain the train was two hours late.
  \item \emph{Recall:} say the Telugu for a railway coach, then the Telugu for a fare, then say \emph{ālasyaṁ} again
\end{itemize}

\begin{tcolorbox}[breakable,colback=teal!4,colframe=teal!35!black,title={Before you move on}]
What does \te{ఆలస్యం} mean? (\textquotedblleft{}A delay; lateness\textquotedblright{}.) Say it once more.
\end{tcolorbox}
\section[\texorpdfstring{\te{వేచి} \te{ఉండు} (vēci uṇḍu)}{వేచి ఉండు (vēci uṇḍu)}]{\te{వేచి} \te{ఉండు} (vēci uṇḍu) — to wait}

\label{lesson:TE-C440-veci-undu}



\begin{quote}
Before the new one: say the Telugu for a fare, then the Telugu for a delay.
\end{quote}

\subsection*{You'll want to know: \te{వేచి} \te{ఉండు}}
\textbf{\te{వేచి} \te{ఉండు}} — \emph{vēci uṇḍu} — \textquotedblleft{}to wait\textquotedblright{}.

\begin{quote}
\te{మేము} \te{గంటసేపు} \te{బస్సు} \te{కోసం} \te{వేచి} \te{ఉన్నాం}.
\end{quote}

\begin{quote}
\emph{mēmu gaṇṭasēpu bassu kōsaṁ vēci unnāṁ.}
\end{quote}

\begin{quote}
We waited an hour for the bus.
\end{quote}

Past (I …): \te{వేచి} \te{ఉన్నాను}, \emph{vēci unnānu}.

What you wait for takes \te{కోసం}: \te{బస్సు} \te{కోసం} \te{వేచి} \te{ఉన్నాను}. It is a little more careful than \te{ఆగు}, to stop and wait.

\begin{cousinweb}
\textbf{\te{వేచి} \te{ఉండు}} (\emph{vēci uṇḍu}) is native Dravidian, a phrase of two verbs: \te{వేచి}, from \te{వేచు}, \textquotedblleft{}to wait for, to look out for\textquotedblright{} (DEDR 5533), and \te{ఉండు}, \textquotedblleft{}to be, to stay\textquotedblright{} (DEDR 697). Together they say \textquotedblleft{}to stay waiting\textquotedblright{}.
\end{cousinweb}

\begin{grammarlens}[title={What you've built}]
One more word for getting around.
\end{grammarlens}

\subsection*{Your turn}
\begin{itemize}
\raggedright
  \item \emph{Say it:} \emph{vēci uṇḍu}
  \item \emph{Say it:} \emph{vēci uṇḍu}, once more
  \item \emph{Say it:} \emph{mēmu gaṇṭasēpu bassu kōsaṁ vēci unnāṁ.} — We waited an hour for the bus.
  \item \emph{Recall:} say the Telugu for a fare, then the Telugu for a delay, then say \emph{vēci uṇḍu} again
\end{itemize}

\begin{tcolorbox}[breakable,colback=teal!4,colframe=teal!35!black,title={Before you move on}]
What does \te{వేచి} \te{ఉండు} mean? (\textquotedblleft{}To wait\textquotedblright{}.) Say it once more.
\end{tcolorbox}
\section[\texorpdfstring{\te{సమయపాలన} (samayapālana)}{సమయపాలన (samayapālana)}]{\te{సమయపాలన} (samayapālana) — punctuality}

\label{lesson:TE-C440-samayapalana}



\begin{quote}
Before the new one: say the Telugu for a delay, then the Telugu for to wait.
\end{quote}

\subsection*{You'll want to know: \te{సమయపాలన}}
\textbf{\te{సమయపాలన}} — \emph{samayapālana} — \textquotedblleft{}punctuality\textquotedblright{}.

\begin{quote}
\te{మా} \te{డ్రైవరుకు} \te{సమయపాలన} \te{చాలా} \te{ముఖ్యం}.
\end{quote}

\begin{quote}
\emph{mā ḍraivaruku samayapālana cālā mukhyaṁ.}
\end{quote}

\begin{quote}
Punctuality matters a great deal to our driver.
\end{quote}

\te{సమయం}, time, + \te{పాలన}, keeping: keeping time. \te{పాలన} alone is rule or administration.

\begin{cousinweb}
\textbf{\te{సమయపాలన}} (\emph{samayapālana}) is a compound of two words borrowed from Sanskrit: \emph{samaya}, \textquotedblleft{}an appointed time\textquotedblright{} (CDIAL 13185), and \emph{pālana}, \textquotedblleft{}guarding, keeping\textquotedblright{} (CDIAL 8128).
\end{cousinweb}

\begin{grammarlens}[title={What you've built}]
One more word for getting around.
\end{grammarlens}

\subsection*{Your turn}
\begin{itemize}
\raggedright
  \item \emph{Say it:} \emph{samayapālana}
  \item \emph{Say it:} \emph{samayapālana}, once more
  \item \emph{Say it:} \emph{mā ḍraivaruku samayapālana cālā mukhyaṁ.} — Punctuality matters a great deal to our driver.
  \item \emph{Recall:} say the Telugu for a delay, then the Telugu for to wait, then say \emph{samayapālana} again
\end{itemize}

\begin{tcolorbox}[breakable,colback=teal!4,colframe=teal!35!black,title={Before you move on}]
What does \te{సమయపాలన} mean? (\textquotedblleft{}Punctuality\textquotedblright{}.) Say it once more.
\end{tcolorbox}
\section[\texorpdfstring{\te{విరామం} (virāmaṁ)}{విరామం (virāmaṁ)}]{\te{విరామం} (virāmaṁ) — an interval; a break}

\label{lesson:TE-C440-viramam}



\begin{quote}
Before the new one: say the Telugu for to wait, then the Telugu for punctuality.
\end{quote}

\subsection*{You'll want to know: \te{విరామం}}
\textbf{\te{విరామం}} — \emph{virāmaṁ} — \textquotedblleft{}an interval; a break\textquotedblright{}.

\begin{quote}
\te{రెండు} \te{గంటల} \te{ప్రయాణం} \te{తరువాత} \te{పది} \te{నిమిషాల} \te{విరామం} \te{ఉంది}.
\end{quote}

\begin{quote}
\emph{reṇḍu gaṇṭala prayāṇaṁ taruvāta padi nimiṣāla virāmaṁ undi.}
\end{quote}

\begin{quote}
After two hours on the road there is a ten-minute break.
\end{quote}

A film has an interval, \te{విరామం}, too. In Sanskrit grammar \emph{virāma} is also the name of the sign \te{్}, which takes away a consonant's vowel.

\begin{cousinweb}
\textbf{\te{విరామం}} (\emph{virāmaṁ}) is borrowed from Sanskrit \emph{virāma}, \textquotedblleft{}cessation, a stop, a pause\textquotedblright{} (Monier-Williams), from \emph{vi-ram}, \textquotedblleft{}to stop\textquotedblright{}.
\end{cousinweb}

\begin{grammarlens}[title={What you've built}]
One more word for getting around.
\end{grammarlens}

\subsection*{Your turn}
\begin{itemize}
\raggedright
  \item \emph{Say it:} \emph{virāmaṁ}
  \item \emph{Say it:} \emph{virāmaṁ}, once more
  \item \emph{Say it:} \emph{reṇḍu gaṇṭala prayāṇaṁ taruvāta padi nimiṣāla virāmaṁ undi.} — After two hours on the road there is a ten-minute break.
  \item \emph{Recall:} say the Telugu for to wait, then the Telugu for punctuality, then say \emph{virāmaṁ} again
\end{itemize}

\begin{tcolorbox}[breakable,colback=teal!4,colframe=teal!35!black,title={Before you move on}]
What does \te{విరామం} mean? (\textquotedblleft{}An interval; a break\textquotedblright{}.) Say it once more.
\end{tcolorbox}
\section[\texorpdfstring{\te{త్వరపడు} (tvarapaḍu)}{త్వరపడు (tvarapaḍu)}]{\te{త్వరపడు} (tvarapaḍu) — to hurry up}

\label{lesson:TE-C440-tvarapadu}



\begin{quote}
Before the new one: say the Telugu for punctuality, then the Telugu for an interval.
\end{quote}

\subsection*{You'll want to know: \te{త్వరపడు}}
\textbf{\te{త్వరపడు}} — \emph{tvarapaḍu} — \textquotedblleft{}to hurry up\textquotedblright{}.

\begin{quote}
\te{బస్సు} \te{వస్తోందని} \te{మేము} \te{త్వరపడ్డాం}.
\end{quote}

\begin{quote}
\emph{bassu vastōndani mēmu tvarapaḍḍāṁ.}
\end{quote}

\begin{quote}
Hearing that the bus was coming, we hurried.
\end{quote}

Past (I …): \te{త్వరపడ్డాను}, \emph{tvarapaḍḍānu}.

\te{త్వరగా}, quickly, holds the same Sanskrit word. \te{తొందరపడు} is its everyday native twin.

\begin{cousinweb}
\textbf{\te{త్వరపడు}} (\emph{tvarapaḍu}) is a Telugu compound: Sanskrit \emph{tvarā}, \textquotedblleft{}haste\textquotedblright{} (Monier-Williams), with the native verb \te{పడు}, \textquotedblleft{}to fall into, to undergo\textquotedblright{} (DEDR 3853).
\end{cousinweb}

\begin{grammarlens}[title={What you've built}]
One more word for getting around.
\end{grammarlens}

\subsection*{Your turn}
\begin{itemize}
\raggedright
  \item \emph{Say it:} \emph{tvarapaḍu}
  \item \emph{Say it:} \emph{tvarapaḍu}, once more
  \item \emph{Say it:} \emph{bassu vastōndani mēmu tvarapaḍḍāṁ.} — Hearing that the bus was coming, we hurried.
  \item \emph{Recall:} say the Telugu for punctuality, then the Telugu for an interval, then say \emph{tvarapaḍu} again
\end{itemize}

\begin{tcolorbox}[breakable,colback=teal!4,colframe=teal!35!black,title={Before you move on}]
What does \te{త్వరపడు} mean? (\textquotedblleft{}To hurry up\textquotedblright{}.) Say it once more.
\end{tcolorbox}
